\section{Setup}
\label{sec:setup}

On-policy distillation (OPD) trains a student on its own samples. The student generates a continuation from a prompt. At every prefix of that continuation, the training loss pulls the student's next-token distribution toward the next-token distribution of a fixed teacher \citep{gu2023minillm,agarwal2024onpolicydistillation,tinker_opd_2025}. This section fixes the notation for the objects that recur in the paper, including the policies, the source distribution that a run trains on, the request field that the teacher induces on that source, and the evaluation targets that we score runs against.

\paragraph{Policies and rollouts.}
A prefix $s$ is a prompt followed by a partial continuation. Prefixes range over a space $\cS$, and tokens range over a finite vocabulary $\cV$. The student has parameters $\theta \in \R^d$, logits $z_\theta(s) \in \R^{|\cV|}$, and next-token distribution $p_\theta(a \mid s) = \softmax(z_\theta(s))_a$. The teacher $q(\cdot \mid s)$ is fixed throughout training. Training proceeds in steps $t = 0, 1, 2, \ldots$. Gradient flow, the continuous-time version of the training process, has the update rule $\dot\theta_t = u_t$. For discrete gradient steps this gives a first-order local approximation \cref{app:ftc}. Two distributions share the student's weights. The \emph{model policy} $p_t$ is the distribution that the loss acts on. The \emph{rollout policy} $\pi_t = \cG[p_t]$ is the distribution that generates the samples (the transform $\cG$ collects the temperature, end-of-sequence conventions, etc). We sample prompts $b\sim\beta$ and continue them under $\pi_t$. The step-$t$ \emph{rollout population} is $\Pi_t(\xi)=\beta(b)\prod_{j}\pi_t(a_j\mid s_j)$ for a trajectory $\xi=(b,a_0,a_1,\ldots)$, with $s_0=b$ and $s_{j+1}=s_j\oplus a_j$ up to termination or the cap. The \emph{training source distribution} $\mu_t = \cM(\Pi_t)$ is the distribution over prefixes that the model trains on ($\cM$ normalizes over the rollouts to form the probability distribution).

\paragraph{Loss, request field, and update.}
\label{def:features}
For reverse Kullback--Leibler divergence, the per-prefix loss is $c_t(s)=\KL{p_t(\cdot\mid s)}{q(\cdot\mid s)}$. Write $\ell_t(s,a) = \log\big(p_t(a \mid s)/q(a \mid s)\big)$ for the log-ratio of the two policies at a token, and $\bar\ell_t(s) = \Ex_{a \sim p_t}[\ell_t(s,a)]$ for its mean under the student. The negative gradient of the loss with respect to the logits is
$B_t(s,a) \;=\; p_t(a \mid s)\,\big[\bar{\ell}_t(s) - \ell_t(s,a)\big]$.
We call $B_t$ the \emph{request field}\footnote{This is the ``student-teacher disagreement'' abstraction our theory is primarily concerned with: we adopt the specific reverse KL due to its popularity in OPD, but the theory covers general disagreement-measuring losses.}. It captures the student-teacher disagreement on a given prefix. The product $\Gamma_t = \mu_t B_t$, visitation times request, is the \emph{request population}. The student's \emph{tangent map} $\Phi_t(s) = \nabla_\theta z_{\theta_t}(s) \in \R^{|\cV| \times d}$. The map pulls a request in logit space back into parameter space, giving the per-prefix gradient $g_t(s) = \Phi_t(s)^{\!\top} B_t(s) = -\nabla_\theta c_t(s)$ and the update direction $u_t = \Ex_{s \sim \mu_t}[g_t(s)]$. See \cref{app:standing} for further details on the setup.
